% Capitulo: Testes e qualidade de dados
\chapter{Testes e qualidade de dados}
\label{cap:testes}

\rascunho{
  \textbf{Ideia do capítulo.} Como saber que o dado está certo antes que o usuário descubra que está errado.
  \begin{itemize}
    \item Tipos de teste: unicidade, nulos, integridade referencial, regras de negócio.
    \item TDD em dados: com IA, escrever testes ficou tão barato que o desenvolvimento
      guiado por testes deixa de ser um \emph{nice to have} e pode virar requisito.
    \item Contratos de dados entre quem produz e quem consome.
    \item Exemplo de teste em SQL.
    \item Relato: um erro em produção que um teste simples teria pego.
  \end{itemize}
}

\section{O que testar}
\label{sec:o-que-testar}

\pendente{escrever}

\section{Testes antes do modelo}
\label{sec:tdd}

\rascunho{
  \textbf{Argumento da seção.} TDD sempre foi recomendado em dados e quase nunca
  praticado, porque escrever teste custava tempo que o prazo não dava. Com IA, o custo
  de escrever o teste caiu, e o TDD pode passar de boa intenção a requisito do time.
  \begin{itemize}
    \item O fluxo: primeiro se escrevem os testes que descrevem as regras de negócio
      (\enquote{toda venda cancelada sai da receita}, \enquote{o saldo nunca fica
      negativo}), depois o modelo que faz os testes passarem.
    \item Os testes viram a especificação executável. A regra é validada uma vez com o
      stakeholder, na forma de casos de teste, e daí em diante o time itera sem
      precisar dele presente a cada versão do modelo.
    \item Consequência para o ritmo: menos dependência da agenda do stakeholder,
      ciclos mais curtos, mudança de regra vira mudança de teste.
    \item O papel da IA: gerar casos de teste a partir da regra descrita em texto,
      sugerir casos de borda, montar dados de exemplo. O papel do profissional:
      revisar se o teste diz mesmo o que o negócio quer (ver o capítulo~\ref{cap:negocio}).
    \item Cuidado: teste gerado sem revisão só confirma o que o modelo já faz, inclusive
      o erro.
    \item Exemplo: escrever os testes de um modelo de receita antes do modelo, em SQL
      ou no formato de testes da ferramenta usada. \pendente{escolher a ferramenta do exemplo}
  \end{itemize}
}

\pendente{escrever}

\section{Contratos de dados}
\label{sec:contratos}

\pendente{escrever}

\section{Um teste em SQL}
\label{sec:teste-sql}

\pendente{escrever}
